% !TeX root = Chapter_07_Slides.tex
% Chapter 7 lecture slides — Loss Control and Risk Reduction
% Instructor-resource series; modeled on the Chapter 1 pilot deck.
\documentclass[11pt, aspectratio=169]{beamer}
\usepackage{tikz}
\makeatletter
\def\input@path{{./}{../slides/}}
\makeatother
\usepackage{erm_beamer_theme}
\usepackage{amsmath,amssymb}

\title[ERM Ch.\,7]{Loss Control and Risk Reduction}
\subtitle{Enterprise Risk Management, Second Edition --- Chapter 7}
\author{Martin F.\ Grace}
\institute{University of Iowa\\[2pt] Tippie College of Business\\[2pt] Vaughan Institute of Risk Management}
\date{June 2026}

\begin{document}

\begin{frame}[plain]
\titlepage
\end{frame}

\begin{frame}{Learning Objectives}
\begin{itemize}
  \item Explain the \alert{frequency--severity} framework and how loss control targets each component.
  \item Calculate the ROI of loss control investments from pre- and post-control frequency and severity data.
  \item Apply the NIOSH \alert{hierarchy of controls} (elimination $\rightarrow$ substitution $\rightarrow$ engineering $\rightarrow$ administrative $\rightarrow$ PPE) to workplace hazards.
  \item Measure loss control effectiveness with TRIR, DART, and EMR.
  \item Evaluate the trade-off between risk financing (insurance, retention) and risk reduction using total cost of risk.
  \item Integrate loss control with the appetite and portfolio frameworks of Chapters 5--6.
\end{itemize}
\end{frame}

\section{Opening Case: Crestline Mutual}

\begin{frame}{The House That Didn't Burn}
\begin{itemize}
  \item Electrical fires in single-family homes: a stubborn line in Crestline Mutual's loss reports --- not the most frequent cause, but disproportionately \emph{severe}.
  \item Traditional responses: higher rates for older homes, smoke-alarm discounts, occasional inspections.
  \item A vendor's different pitch: a night-light-sized device that ``listens'' to the home's wiring, detecting arcing and failing connections \emph{before} visible signs appear.
  \item The actuaries asked for data: in a multi-carrier pilot across 1M+ home-years, device homes had measurably fewer electrical-fire claims than a matched control group.
\end{itemize}
\end{frame}

\begin{frame}{What the Data Showed}
\begin{columns}[T]
\column{0.48\textwidth}
\textbf{One house}
\begin{itemize}
  \item October alert: repeated arcing on a garage circuit --- no flickering lights, no tripped breakers
  \item Electrician finds a charred, back-stabbed outlet terminal: ``the kind of thing that smolders before it lights something up''
  \item Repair: under an hour, a screwdriver, and a new outlet
\end{itemize}
\column{0.48\textwidth}
\textbf{Thousands of houses}
\begin{itemize}
  \item Participating homes vs.\ control group: fewer electrical-fire claims (\alert{frequency})
  \item Of fires that still occurred, more contained to a room or appliance (\alert{severity})
  \item The \emph{shape of the loss distribution} changed --- not the policy form, the underlying risk
\end{itemize}
\end{columns}
\end{frame}

\begin{frame}{The Lesson}
\begin{itemize}
  \item Loss control changes the physical reality that generates losses --- it doesn't reprice or transfer them.
  \item Crestline ran the chapter's core arithmetic across a portfolio: device and monitoring cost per home vs.\ reduced frequency, reduced severity, and avoided adjustment expense.
  \item Its advantage: a portfolio large enough to make the reductions statistically visible. A single firm must lean on industry data and vendor studies instead.
\end{itemize}
\medskip
\begin{block}{The decision problem}
You run the actuarial team. The traditional toolkit --- price the risk, adjust the deductible --- is well understood. Prevention requires spending money on \emph{other people's houses} and waiting years for proof. What evidence would you require before funding it?
\end{block}
\end{frame}

\section{From Identification to Reduction}

\begin{frame}{The Fourth Response: Reduce}
\begin{itemize}
  \item The risk response menu: \textbf{avoid}, \textbf{reduce}, \textbf{transfer}, \textbf{accept}. Avoidance forgoes profit; transfer shifts financial consequences at a price; acceptance is for risks within appetite. This chapter is about \alert{reduction}.
\end{itemize}
\medskip
\begin{block}{Loss control}
All activities designed to reduce the \alert{frequency} (likelihood) of loss events or the \alert{severity} (magnitude) of losses when events occur --- changing the underlying physical and operational reality.
\end{block}
\medskip
\begin{itemize}
  \item Side benefits compound: lower premiums via experience rating, better efficiency, morale and retention, compliance, less litigation, reputation.
  \item ERM tie-in: risks outside appetite demand controls regardless of ROI; risks that correlate with other exposures (Ch.~6) deserve priority --- preventing one factory fire avoids property, interruption, contract, and reputational losses at once.
\end{itemize}
\end{frame}

\section{The Economics of Loss Control}

\begin{frame}{Total Cost of Risk}
\[
\begin{aligned}
\text{TCOR} ={}& \text{Direct Losses} + \text{Loss Control Costs} + \text{Risk Transfer Costs} \\
&+ \text{Administrative Costs} + \text{Indirect Costs}
\end{aligned}
\]
\begin{itemize}
  \item Typically 1--5\% of revenue; higher in construction, transportation, manufacturing.
  \item Loss control appears as a \emph{cost} --- but it simultaneously cuts direct losses, transfer costs (better experience rating), and indirect costs (fewer disruptions).
\end{itemize}
\medskip
\begin{block}{The objective}
Minimize \emph{total} cost of risk, not any single component --- the answer when a CFO asks why the safety budget is growing.
\end{block}
\end{frame}

\begin{frame}{Loss Control ROI}
\[
\text{Expected Loss} = \text{Frequency} \times \text{Average Severity}
\]
\[
\text{ROI} = \frac{\text{Annual Benefit} - \text{Annualized Control Cost}}{\text{Annualized Control Cost}}
\qquad
\text{Payback} = \frac{\text{Initial Investment}}{\text{Annual Benefit}}
\]
\begin{exampleblock}{Example: forklift--pedestrian collisions in a warehouse}
\begin{itemize}
  \item Pre-control: $8 \times \$12{,}000 = \$96{,}000$/yr. Controls (barriers, markings, proximity warnings): \$45{,}000 capital + \$3{,}000/yr training
  \item Post-control: frequency $-75\%$, severity $-30\%$ $\Rightarrow$ $2 \times \$8{,}400 = \$16{,}800$/yr
  \item Annualized cost \$7{,}500; savings \$79{,}200 $\Rightarrow$ \alert{ROI = 956\%, payback $\approx$ 7 months}
\end{itemize}
\end{exampleblock}
\medskip
Returns like this would be extraordinary anywhere else in the capital budget. In loss control, they are not unusual --- though marginal returns diminish as the obvious hazards get fixed.
\end{frame}

\section{Frequency vs.\ Severity}

\begin{frame}{Frequency Reduction: Fewer Events}
\begin{itemize}
  \item Shifts the loss \emph{count} distribution leftward; most valuable for high-frequency, lower-severity risks.
  \item \textbf{Engineering}: machine guarding (29 CFR 1910.212), interlocks and fail-safes, automatic systems, physical separation of pedestrians from forklifts.
  \item \textbf{Administrative}: lockout/tagout (OSHA estimates $\sim$120 fatalities and 50{,}000 injuries prevented annually), hot-work and confined-space permits, job safety analysis, inspections and near-miss reporting.
  \item \textbf{Training and culture}: hands-on competency, refreshers, leadership that treats safety as a value rather than compliance theater.
  \item Crestline's alerts worked here: incipient fires became electrician visits --- events removed from the loss count entirely.
\end{itemize}
\end{frame}

\begin{frame}{Severity Control: Smaller Losses When Events Occur}
\begin{itemize}
  \item Accepts that some events happen; compresses the severity distribution --- average loss \emph{and} the tail.
  \item \textbf{Property}: sprinklers cut average loss per fire 60--70\% and civilian fire deaths 87\% (NFPA); compartmentation; separating high-value assets across locations; secondary containment.
  \item \textbf{Injury}: emergency response and AEDs; early return-to-work (claim durations $-30$ to $-50\%$); active claims management ($-20$ to $-30\%$ medical cost per claim); PPE as the last layer.
  \item \textbf{Business interruption}: redundancy, tested continuity plans (2--3$\times$ faster recovery), supply chain diversification. Maersk's post-NotPetya segmentation is severity control --- it wouldn't prevent the next infection, but it would \emph{contain} it.
\end{itemize}
\end{frame}

\section{The Hierarchy of Controls}

\begin{frame}{The NIOSH Hierarchy}
\begin{block}{Most effective $\rightarrow$ least effective}
\alert{Elimination} $\rightarrow$ \alert{Substitution} $\rightarrow$ \alert{Engineering} $\rightarrow$ \alert{Administrative} $\rightarrow$ \alert{PPE}
\end{block}
\medskip
\begin{itemize}
  \item \textbf{Elimination}: remove the hazard (reformulate paint to drop hexavalent chromium). Absolute --- and often infeasible for core processes.
  \item \textbf{Substitution}: water-based solvents for VOC solvents. Check the replacement doesn't import new hazards.
  \item \textbf{Engineering}: guards, ventilation, automation --- works passively, no behavior change required.
  \item \textbf{Administrative}: rotation, LOTO, permits --- depends on consistent human compliance.
  \item \textbf{PPE}: last line of defense; doesn't reduce the hazard, depends entirely on proper individual use.
\end{itemize}
\medskip
The ranking is about \emph{reliability}: how vulnerable is the control to human error?
\end{frame}

\begin{frame}{Applying the Hierarchy: A 95 dBA Stamping Press}
\begin{itemize}
  \item \textbf{Elimination}: stamping is essential --- infeasible.
  \item \textbf{Substitution}: quieter hydraulic presses at the next equipment replacement cycle?
  \item \textbf{Engineering}: sound enclosure, vibration damping, relocation --- below 85 dBA, the hazard is gone for most workers.
  \item \textbf{Administrative}: job rotation to cap cumulative dose; schedule noisy work for low-occupancy hours.
  \item \textbf{PPE}: hearing protection for the residual --- and as redundant backup even with engineering controls.
\end{itemize}
\medskip
\begin{block}{Defense-in-depth}
Real programs combine levels. Economics align with effectiveness: higher-hierarchy controls cost more upfront but less over the lifecycle --- and OSHA requires feasible engineering controls \emph{before} reliance on PPE.
\end{block}
\end{frame}

\section{Loss Control Across Industries}

\begin{frame}{Worked Example: Machine Guarding at Mercer Tool \& Stamping}
\begin{exampleblock}{150-employee metal fabricator; 18 machine injuries in 3 years, \$627{,}000 total cost (\$209{,}000/yr)}
\begin{itemize}
  \item \$125{,}000 program: interlocked guards, light curtains, fixed guarding, training, inspections
  \item Frequency: 6/yr $\rightarrow$ 1.2/yr ($-80\%$, conservative); severity: \$34{,}833 $\rightarrow$ \$20{,}900 ($-40\%$)
  \item Post-control loss: $1.2 \times \$20{,}900 = \$25{,}080$; annual savings \$183{,}920
  \item Annualized cost \$20{,}500 $\Rightarrow$ \alert{ROI = 797\%, payback $\approx$ 8 months}
\end{itemize}
\end{exampleblock}
\medskip
Three years on: one minor injury vs.\ 3.6 predicted; EMR 1.24 $\rightarrow$ 0.89, saving \$67{,}000/yr in premium --- a benefit not even counted in the original analysis.
\end{frame}

\begin{frame}{Four Industries, One Template}
Baseline frequency and severity $\rightarrow$ design controls $\rightarrow$ post-control profile $\rightarrow$ compute the return.
\medskip
\begin{center}
\small
\begin{tabular}{llrrr}
\toprule
\textbf{Case} & \textbf{Hazard} & \textbf{Investment} & \textbf{ROI} & \textbf{Payback} \\
\midrule
Mercer Tool (mfg.) & Machine contact & \$125K & 797\% & 8 mo. \\
Midtown Foods (retail) & Slip-and-fall & \$180K & 780\% & 6 mo. \\
Summit Builders (constr.) & Falls from elevation & \$340K & 428\% & 9 mo. \\
Community Regional (health) & Workplace violence & \$425K & 103\% & 2.3 yr. \\
\bottomrule
\end{tabular}
\end{center}
\medskip
\begin{itemize}
  \item Midtown's grape decision is \emph{elimination}; its matting is engineering; its checklists are administrative --- the hierarchy in a grocery store.
  \item Even the weakest result (103\%) would be celebrated in most capital budgets --- and it excludes retention, reputation, and a safety-climate score that rose from 61\% to 84\%.
\end{itemize}
\end{frame}

\section{Measuring Effectiveness}

\begin{frame}{Lagging Indicators: TRIR and DART}
\[
\text{TRIR} = \frac{\text{Recordable Cases} \times 200{,}000}{\text{Total Hours Worked}}
\]
\begin{itemize}
  \item The 200{,}000 constant normalizes to 100 full-time workers --- rates comparable across firm sizes.
  \item \alert{DART} counts the serious subset: days away, restricted duty, or transfer. High DART relative to TRIR signals injuries trending severe.
  \item 2022 benchmarks: all private industry TRIR 2.7; manufacturing 3.2; healthcare 4.5; construction 2.4.
\end{itemize}
\begin{exampleblock}{Example}
450 employees, 900{,}000 hours, 18 recordables: TRIR $= (18 \times 200{,}000)/900{,}000 = 4.0$ --- above the 3.2 manufacturing average. Investigate.
\end{exampleblock}
\end{frame}

\begin{frame}{EMR: Where Safety Meets the Premium}
\begin{itemize}
  \item \alert{Experience Modification Rate}: actual losses vs.\ expected losses for similar firms. EMR = 1.00 is average; below 1.00 earns discounts, above 1.00 pays surcharges.
\end{itemize}
\[
\text{Premium} = \text{Manual Premium} \times \text{EMR} \times \text{Other Modifiers}
\]
\begin{exampleblock}{Example}
Manual premium \$500{,}000; EMR falls 1.25 $\rightarrow$ 0.95 through loss control: savings $= \$500{,}000 \times 0.30 = \alert{\$150{,}000}$ per year.
\end{exampleblock}
\medskip
\begin{itemize}
  \item Timing matters: EMR uses a three-year experience window with a one-year lag --- premium benefits arrive in years 4--7, not year 1.
  \item Frequency is weighted heavily in the formula: even small claims hurt the mod.
\end{itemize}
\end{frame}

\begin{frame}{Leading vs.\ Lagging Indicators}
\begin{columns}[T]
\column{0.48\textwidth}
\textbf{Lagging (outcomes)}
\begin{itemize}
  \item TRIR, DART, lost workdays, EMR, claim costs
  \item Necessary --- but backward-looking, and sparse data when injuries are rare
  \item 18 injury-free months can mean safety, or accumulating undetected hazards
\end{itemize}
\column{0.48\textwidth}
\textbf{Leading (proactive activities)}
\begin{itemize}
  \item Training currency, inspections completed, hazard-correction closure time
  \item Near-miss reports (typically 10--30 per recordable --- a far richer dataset)
  \item Observation counts, pre-task planning compliance, preventive maintenance on time
\end{itemize}
\end{columns}
\medskip
Crestline again: anomalies flagged and faults fixed were the \emph{leading} indicators; fire claims, months later, the lagging confirmation. Mature programs run a balanced scorecard of both.
\end{frame}

\section{Risk Financing vs.\ Risk Reduction}

\begin{frame}{Control First, Then Retain}
A retail chain weighs a deductible increase: \$25K $\rightarrow$ \$100K per claim, saving \$225{,}000 in premium. 32 claims/yr.
\medskip
\begin{center}
\small
\begin{tabular}{lrrrr}
\toprule
\textbf{Strategy} & \textbf{Premium} & \textbf{Retained} & \textbf{Control} & \textbf{Total} \\
\midrule
\$25K deductible, no control & \$850K & \$576K & --- & \$1{,}426K \\
\$100K deductible alone & \$625K & \$1{,}156K & --- & \$1{,}781K \\
\$100K deductible + loss control & \$625K & \$470K & \$40K & \alert{\$1{,}135K} \\
\bottomrule
\end{tabular}
\end{center}
\medskip
\begin{itemize}
  \item The deductible increase alone \emph{costs} \$355K/yr; bundled with slip-and-fall controls it \emph{saves} \$291K/yr.
  \item Sequencing is the lesson: the higher retention becomes safe only after the underlying risk is reduced. Same logic drives the strong control incentives in self-insurance and captives --- there's no insurer to pass costs to.
\end{itemize}
\end{frame}

\section{Capstone: Riverside Metal Products}

\begin{frame}{Riverside: Reshaping the Whole Distribution}
150-employee fabricator; three-year program targeting strains, lacerations, and slips.
\medskip
\begin{center}
\small
\begin{tabular}{lrr}
\toprule
 & \textbf{Pre-control} & \textbf{Post-control} \\
\midrule
Claims per year & 12.5 & 4.2 \\
Average claim cost & \$12{,}400 & \$9{,}800 \\
Expected annual loss & \$155{,}000 & \$41{,}160 \\
95\% VaR & \$225{,}000 & \$78{,}000 \\
TRIR (industry avg.\ 3.2) & 8.3 & 2.8 \\
EMR & 1.42 & 0.87 \\
\bottomrule
\end{tabular}
\end{center}
\medskip
Expected loss fell 73\% --- but the 95\% VaR fell 65\% too. The program compressed the \emph{tail}, not just the mean: the Crestline signature inside a single firm.
\end{frame}

\begin{frame}{Riverside: The Financial Case}
\begin{exampleblock}{Investment: \$230{,}000 capital over 3 years + \$22{,}000/yr ongoing (total annual cost \$45{,}000)}
\begin{itemize}
  \item Direct loss reduction: $\$155{,}000 - \$41{,}160 = \$113{,}840$/yr
  \item Premium savings (EMR 1.42 $\rightarrow$ 0.87): $\$285{,}000 - \$165{,}000 = \$120{,}000$/yr
  \item Total benefit \$233{,}840 $\Rightarrow$ \alert{ROI = 420\%, payback $\approx$ 12 months}; five-year cumulative savings $\approx$ \$829{,}000
\end{itemize}
\end{exampleblock}
\medskip
\begin{itemize}
  \item Premium savings roughly \emph{doubled} the benefit --- experience rating is half the business case.
  \item Twelve-month payback competes with any capital project, yet safety investments routinely face more scrutiny than production investments with weaker returns.
\end{itemize}
\end{frame}

\section{From ROI to NPV (Appendix 7A)}

\begin{frame}{Loss Control as a Capital Investment}
\begin{itemize}
  \item The chapter's ROI model assumes no discounting, no taxes, static premiums. Appendix 7A relaxes each. The NPV criterion:
\end{itemize}
\[
NPV = (\Delta L + \Delta P - C)(1-\tau) \times PVIFA(r,T) - I_0(1-\tau) \;>\; 0
\]
\begin{itemize}
  \item Tax subtlety: losses are deductible, so \emph{preventing} a \$100K loss saves only \$75K after tax at $\tau = 0.25$ --- prevented losses are worth $(1-\tau)$ on the dollar.
  \item Mercer Tool revisited: the 797\% simple ROI corresponds to an after-tax NPV $\approx$ \$792K; adding the (lagged) premium savings raises it to $\approx$ \$896K --- on a net after-tax outlay of \$93{,}750.
  \item Discount rate: loss-prevention cash flows pay off in bad states --- arguments exist for rates below WACC. Using WACC is the conservative test.
\end{itemize}
\end{frame}

\begin{frame}{When Expected Value Isn't Enough}
\begin{exampleblock}{Example: a chemical plant's \$15M containment upgrade}
\begin{itemize}
  \item MFL \$500M, frequency 1/200 years; upgrade cuts MFL to \$150M
  \item Expected-value NPV at 8\%, 10 years: about $-\$3.3$M --- \emph{negative}
  \item But \$200M equity + \$300M debt: a \$500M loss means bankruptcy; \$150M is survivable
  \item Valuing avoided distress costs at \$100M adds $\approx$ \$3.4M --- total NPV turns positive
\end{itemize}
\end{exampleblock}
\medskip
\begin{itemize}
  \item Loss control that compresses the tail creates value through three channels: expected loss, VaR, and the probability of \alert{financial distress} --- the simplified model sees only the first.
  \item Real-options footnote: Crestline waiting for multi-carrier pilot data before launching was deferral-option logic --- the pilot resolved effectiveness uncertainty cheaply.
\end{itemize}
\end{frame}

\begin{frame}{Discussion Questions}
\begin{enumerate}
  \item Distinguish frequency reduction from severity control with an example of each from different industries. When would you invest in severity control rather than attempting frequency reduction?
  \item A shear is activated by foot pedal while workers' hands align the sheet near the blade path. Apply all five levels of the hierarchy of controls; assess feasibility and effectiveness at each level.
  \item A manufacturer has a \$425{,}000 manual workers' compensation premium and EMR of 1.28. Compute the current premium, the premium if loss control brings EMR to 0.92, and the annual savings.
  \item Your budget covers loss control for only some of: cyber, supply chain, product liability, workplace injuries. Using Chapter 6's portfolio logic, how do you prioritize --- and what beyond expected loss reduction belongs in the decision?
\end{enumerate}
\end{frame}

\end{document}
